\section{When to Use Merge-Split}
\label{sec:when-to-use}

\subsection{Decision Framework}

The choice among standard \recom{}, \mergesplit{}, \forestrecom{}, and
SMC~\citep{dellaLibera2026smc} depends on three factors:

\begin{enumerate}
  \item \textbf{Distributional requirement.} Does the application require a
        formal target-distribution proof, or is a declared approximate
        reversible-ReCom correction sufficient?
  \item \textbf{Computational budget.} How many steps (and wall-clock
        time) are available?
  \item \textbf{Ensemble size.} Is the goal to produce a large ensemble
        ($\geq 10{,}000$ plans) or a small focused ensemble
        ($\leq 1{,}000$ plans)?
\end{enumerate}

\subsection{Use Merge-Split When}

\mergesplit{} is the recommended chain when:

\begin{itemize}
  \item A declared approximate correction is sufficient --- the application
        benefits from a reproducible reverse-tree ratio and visible step
        diagnostics, but does not require formal exact-determinant guarantees.
        This covers most expert-witness and legislative-disclosure contexts
        where the practitioner can state the actual mechanism used by the
        chain and its approximation boundary.
        We note that no published redistricting court opinion has
        specifically addressed the distinction between exact and approximate
        uniform targeting in Markov chain methods; this claim reflects
        practitioner consensus rather than established legal precedent.
        Consult legal counsel for litigation contexts.

  \item Speed and deterministic candidate selection matter --- the run must
        visit many states and return a requested edge-cut percentile under a
        fixed seed schedule.

  \item Standard ReCom's distributional ambiguity is a concern --- when
        opposing counsel or a court scrutinises the ensemble methodology,
        the unknown stationary distribution of standard \recom{} is a
        vulnerability.
        \mergesplit{} provides a principled, citable response~\citep{janson2022},
        provided the report also states the pair-count and determinant terms
        that are outside the current implementation.

  \item The acceptance rate of \forestrecom{} ($\approx\!40\%$) is too
        low for the available computational budget --- \mergesplit{}'s
        $\approx\!60\%$ acceptance rate produces~50\% more accepted plans
        per step, which may be decisive for small-budget applications.
        Acceptance rate for Forest \recom{} is from our G.9 study
        \citep{dellaLibera2026forestRecom}, which reports 41--48\% on
        NC/WI/TX.
\end{itemize}

\subsection{Use Forest ReCom When}

\forestrecom{}~\citep{dellaLibera2026forestRecom} is preferred when:

\begin{itemize}
  \item Stronger theoretical guarantees are required --- the application must
        demonstrate a forest-measure or determinant-backed target model rather
        than the \mergesplit{} two-tree proxy.
        This matters for peer-reviewed publication claims about exact target
        distributions.

  \item The merged regions are small enough that richer forest-count accounting
        is affordable and useful.

  \item The ensemble is small and variance in the ratio estimate is a
        concern --- with fewer than 100 steps, the random estimate in
        \mergesplit{} may introduce systematic bias in the seat
        distribution, whereas a stronger \forestrecom{} correction is less
        exposed to this specific one-sample variance.
\end{itemize}

\subsection{Use SMC When}

Sequential Monte Carlo~\citep{dellaLibera2026smc} is the right escalation when
the desired claim is a calibrated particle estimate for a declared target:

\begin{itemize}
  \item Stronger calibration is required --- SMC produces weighted particles
        for a declared target and reports effective sample size, resampling,
        and weight degeneracy diagnostics rather than relying on MCMC mixing
        alone.

  \item The ensemble must make its proposal and reweighting stages explicit,
        because the particle filter gives a different audit surface than a
        single Markov chain.

  \item The chain mixing time is a concern --- SMC replaces one long Markov
        path with staged proposal, weighting, and resampling diagnostics.
\end{itemize}
The price of SMC is implementation complexity and memory: retaining and
resampling a particle population requires $O(N \cdot n)$ memory for $N$
particles.
For production 50-state runs, \mergesplit{} is usually preferred.

\subsection{Practical Guidance}

For standard 50-state redistricting analysis:

\begin{enumerate}
  \item Generate the initial plan using \cs{}~\citep{b11paper} or
        another global search method.
  \item Run \mergesplit{} for 1{,}000--10{,}000 steps to visit candidate
        states.
  \item For litigation contexts, state explicitly that the chain uses the
        Merge-Split proposal~\citep{janson2022} with a two-tree stochastic
        correction, reverse zero-cut rejection, and no pair-count multiplier.
  \item If the state has large districts ($n/k \geq 200$ tracts per
        district), prefer a longer run to compensate for the larger
        merged-region size.
  \item Record the base seed, step count, percentile, chosen plan hash, and
        acceptance/cut-count diagnostics for reproducibility.
\end{enumerate}

\paragraph{CLI reference.}
\begin{verbatim}
# 1000 steps; choose minimum-EC candidate
bisect state --state NC --year 2020 \
    --search merge-split --merge-split-steps 1000 --percentile 0.0

# 5000-step chain, still selecting the minimum-EC visited candidate
bisect state --state TX --year 2020 \
    --search merge-split --merge-split-steps 5000 --percentile 0.0
\end{verbatim}
